\section{Experiments}
\label{sec:experiments}

\subsection{Experimental Setup}

\textbf{Datasets.}
We evaluate on three benchmarks:
ETTh1 \cite{zhou2021ett} (17{,}420 time-steps, 7 variates, hourly electricity transformer temperature);
ETTm1 \cite{zhou2021ett} (69{,}680 time-steps, 7 variates, 15-minute resolution);
and Exchange-Rate \cite{wu2021autoformer} (7{,}588 time-steps, 8 currency exchange rates).
For each dataset we report three forecast horizons: $H \in \{96, 336, 720\}$ with a fixed lookback window $L = 96$.
Data are partitioned 70\%, 10\%, and 20\% into training, validation, and test sets, respectively, and normalised per variate with Reversible Instance Normalisation (RevIN \cite{kim2022revin}).
Note that this split differs from the standard chronological split used in published iTransformer and S-Mamba papers; our reported numbers are therefore not directly comparable to those benchmarks.

\textbf{Structures.}
All experiments use two structural templates: S-Mamba (4 genome slots split evenly between variate and temporal backbones) and iTransformer (4 genome slots alternating cross-variate and per-variate sub-layers).

\textbf{Training.}
During NSGA-II search, each candidate genome is evaluated for $N_{\mathrm{eval}} = 500$ steps using the Adam optimiser with learning rate $10^{-3}$.
The top Pareto-front genomes are retrained from scratch for 20{,}000 steps.
All experiments run on a single GPU.

\textbf{Metrics.}
We report test MSE and MAE on un-normalised (RevIN-restored) outputs, consistent with the iTransformer benchmark protocol.

\subsection{Experiment 1: Architecture Search Results}
\label{sec:exp_search}

We run NSGA-II for 18 generations (population 16) on each (dataset, horizon) pair and retrain the Pareto-front genomes for 20{,}000 steps.
Tables~\ref{tab:exp1_etth1} and~\ref{tab:exp1_exchange} report the best-MSE candidate per condition for ETTh1 and Exchange-Rate; ETTm1 follows a similar pattern and is omitted for brevity.

\begin{table}[t]
\caption{Architecture search results on ETTh1. Best test MSE genome per (structure, $H$) reported after 20{,}000-step retraining. Genome entries are operator class IDs from the pool (Table~\ref{tab:operator_pool} and differential variants).}
\label{tab:exp1_etth1}
\centering
\small
\begin{tabular}{llrrr}
\toprule
Structure & $H$ & Params & MSE & MAE \\
\midrule
S-Mamba   & 96  & 528K  & 0.596 & 0.547 \\
S-Mamba   & 336 & 678K  & 0.758 & 0.641 \\
S-Mamba   & 720 & 1{,}590K & 0.819 & 0.670 \\
\midrule
iTransformer & 96  & 1{,}188K & 0.549 & 0.502 \\
iTransformer & 336 & 939K  & 0.809 & 0.661 \\
iTransformer & 720 & 1{,}976K & 0.808 & 0.671 \\
\bottomrule
\end{tabular}
\end{table}

\begin{table}[t]
\caption{Architecture search results on Exchange-Rate. Best test MSE genome per (structure, $H$).}
\label{tab:exp1_exchange}
\centering
\small
\begin{tabular}{llrrr}
\toprule
Structure & $H$ & Params & MSE & MAE \\
\midrule
S-Mamba   & 96  & 570K  & 0.143 & 0.284 \\
S-Mamba   & 336 & 490K  & 0.770 & 0.681 \\
S-Mamba   & 720 & 1{,}059K & 1.084 & 0.829 \\
\midrule
iTransformer & 96  & 1{,}090K & 0.142 & 0.281 \\
iTransformer & 336 & 1{,}153K & 0.673 & 0.629 \\
iTransformer & 720 & 1{,}059K & 1.309 & 0.838 \\
\bottomrule
\end{tabular}
\end{table}

Across all datasets and horizons, the top Pareto-front genomes are consistently \emph{heterogeneous}: no single operator class fills all four slots.
Table~\ref{tab:discovered_genomes} lists the best-discovered genome per condition for ETTh1.
For example, the best ETTh1 H\,=\,96 S-Mamba genome uses all-attention operators across slots — SA-3 (MQA), SA-1 (MHA), SA-3, SA-4 (GQA) — achieving 528K parameters.
The best iTransformer genome for the same condition places GConv-1 (short gated convolution) and GMemless (FFN) across all four slots, a heterogeneous assignment the original iTransformer architecture does not employ.
At longer horizons (H\,=\,720), Rec-5 (a compact CfC variant with 2$\times$ internal expansion) appears in both structures, providing recurrent inductive bias without the parameter cost of full CfC (16$\times$ expansion).
On Exchange-Rate at H\,=\,96, both structures discover compact (${\le}1.1$M parameter) heterogeneous genomes achieving MSE\,$\approx$\,0.14.

\begin{table}[t]
\caption{Best-discovered genomes on ETTh1. S-Mamba slots: 1--2 = Stage-1 (variate), 3--4 = Stage-2 (temporal). iTransformer slots alternate cross-variate (cv) and per-variate (pv). GC-1 = GConv-1; dSA = differential SA; Rec-5 = compact CfC ($2\times$ expansion).}
\label{tab:discovered_genomes}
\centering
\small
\setlength{\tabcolsep}{3pt}
\begin{tabular}{llccccr}
\toprule
Struct. & $H$ & Slot 1 & Slot 2 & Slot 3 & Slot 4 & MSE \\
\midrule
\multirow{3}{*}{S-Mamba}
 & 96  & SA-3    & SA-1    & SA-3   & SA-4     & 0.596 \\
 & 336 & SA-1    & SA-1    & GC-1   & GC-1     & 0.758 \\
 & 720 & GC-1    & SA-4    & Rec-5  & GMemless & 0.819 \\
\midrule
\multirow{3}{*}{iTransf.}
 & 96  & GC-1    & GMemless & GC-1  & GC-1     & 0.549 \\
 & 336 & SA-1    & dSA-1   & SA-4   & dSA-3    & 0.809 \\
 & 720 & GC-1    & Rec-5   & Rec-5  & GMemless & 0.808 \\
\bottomrule
\end{tabular}
\end{table}

\subsection{Experiment 2: Operator Ablation}
\label{sec:exp_ablation}

\begin{figure}[t]
\centering
\includegraphics[width=0.97\columnwidth]{pareto_front.pdf}
\caption{NSGA-II search trajectory for S-Mamba on ETTh1 H\,=\,96. Y-axis is validation MSE (the NSGA-II objective); x-axis is parameter count in millions. Gray dots: all 262 evaluated genomes. Blue line/dots: Pareto-optimal front. Named markers: homogeneous baselines evaluated under the same conditions. The Pareto front sits in the low-parameter, competitive-MSE region that no single-operator baseline occupies.}
\label{fig:pareto}
\end{figure}

Figure~\ref{fig:pareto} shows the full NSGA-II search trajectory and the resulting Pareto front for S-Mamba on ETTh1 H\,=\,96.
The evolutionary search explores 262 genomes; the Pareto front consists of heterogeneous genomes that achieve competitive validation MSE at 0.5M--2M parameters, well below the parameter cost of ALL\_FFN (3.2M) and ALL\_CFC (16.9M).
To isolate the contribution of operator heterogeneity, we compare five baseline configurations against the best NSGA-II genome (MIXED\_GA):
\textbf{ALL\_CONV} (all GConv-1 slots),
\textbf{ALL\_ATTN} (all SA-1 slots),
\textbf{ALL\_FFN} (all GMemless slots),
\textbf{ALL\_CFC} (all CfC slots).
All variants use the same structure template and are trained for 20{,}000 steps.

\begin{table*}[t]
\caption{Operator ablation on ETTh1. Test MSE for each single-class baseline and the best NSGA-II heterogeneous genome (MIXED\_GA). Parameter counts vary with horizon due to the forecast head size.}
\label{tab:exp2_etth1}
\centering
\small
\setlength{\tabcolsep}{4pt}
\begin{tabular}{llrrrrr}
\toprule
 & & \multicolumn{2}{c}{H = 96} & \multicolumn{2}{c}{H = 336} & H = 720 \\
\cmidrule(lr){3-4}\cmidrule(lr){5-6}\cmidrule(lr){7-7}
Structure & Variant & Params & MSE & Params & MSE & MSE \\
\midrule
\multirow{5}{*}{S-Mamba}
 & ALL\_CONV  & 582K  & 0.518 & 644K  & 0.700 & 0.742 \\
 & ALL\_ATTN  & 874K  & 0.548 & 936K  & 0.682 & 0.805 \\
 & ALL\_FFN   & 3.2M  & 0.450 & 3.3M  & 0.549 & 0.657 \\
 & ALL\_CFC   & 16.9M & 0.493 & 17.0M & 0.705 & 0.758 \\
 & MIXED\_GA  & 578K  & 0.537 & 790K  & 0.680 & 0.853 \\
\midrule
\multirow{5}{*}{iTransformer}
 & ALL\_CONV  & 582K  & 0.592 & 644K  & 0.806 & 0.809 \\
 & ALL\_ATTN  & 874K  & 0.555 & 936K  & 0.663 & 0.757 \\
 & ALL\_FFN   & 3.2M  & 0.454 & 3.3M  & 0.551 & 0.668 \\
 & ALL\_CFC   & 16.9M & 0.491 & 17.0M & 0.690 & 0.746 \\
 & MIXED\_GA  & 1.2M  & \textbf{0.464} & 1.1M & 0.677 & 0.749 \\
\bottomrule
\end{tabular}
\end{table*}

\begin{table}[t]
\caption{Operator ablation on Exchange-Rate. Test MSE for single-class baselines vs.\ MIXED\_GA.}
\label{tab:exp2_exchange}
\centering
\small
\setlength{\tabcolsep}{4pt}
\begin{tabular}{llrrr}
\toprule
Structure & Variant & Params & MSE$_{96}$ & MSE$_{336}$ \\
\midrule
\multirow{5}{*}{S-Mamba}
 & ALL\_CONV  & 582K  & 0.672 & 1.372 \\
 & ALL\_ATTN  & 874K  & 0.156 & 0.618 \\
 & ALL\_FFN   & 3.2M  & 0.108 & 0.625 \\
 & ALL\_CFC   & 16.9M & 0.227 & 1.090 \\
 & MIXED\_GA  & 1.6M  & 0.148 & 0.726 \\
\midrule
\multirow{5}{*}{iTransformer}
 & ALL\_CONV  & 582K  & 0.489 & 1.457 \\
 & ALL\_ATTN  & 874K  & 0.203 & 0.763 \\
 & ALL\_FFN   & 3.2M  & \textbf{0.099} & 0.617 \\
 & ALL\_CFC   & 16.9M & 0.393 & 1.289 \\
 & MIXED\_GA  & 1.0M  & 0.158 & 0.795 \\
\bottomrule
\end{tabular}
\end{table}

Tables~\ref{tab:exp2_etth1} and~\ref{tab:exp2_exchange} reveal three consistent patterns.
No single operator dominates across datasets: ALL\_FFN achieves the best ETTh1 MSE but at 3.2M parameters — 2.7--5.5$\times$ more than MIXED\_GA (578K--1.2M); on Exchange-Rate, ALL\_FFN wins on absolute MSE at short horizons while MIXED\_GA matches ALL\_ATTN at one-fifth the parameter cost.
ALL\_CFC's 16.9M parameters (due to $16{\times}$ internal expansion) make it the least efficient variant despite competitive accuracy, confirming that CfC's stability advantage cannot be justified when deployed homogeneously.
The core result holds across both datasets and both templates: NSGA-II finds heterogeneous genomes that occupy Pareto-optimal operating points no single-class baseline reaches.

\subsection{Experiment 3: TiDAR Inference Speedup}
\label{sec:tidar_speedup}

We evaluate TiDAR's inference trade-off between speed and accuracy on Exchange-Rate (H\,=\,96) using the best S-Mamba genome from Experiment~1.
Five operating modes are compared: pure draft ($k=0$), partial-AR with $k \in \{1, 4, 16\}$ autoregressive refinement steps, and full-AR ($k=H=96$).
Latency is measured as median milliseconds per sample over 86 batches.

\begin{table}[t]
\caption{TiDAR inference speed-accuracy trade-off on Exchange-Rate H\,=\,96, S-Mamba. Speedup is relative to full-AR ($k=96$). MSE is on the test set.}
\label{tab:exp4}
\centering
\small
\begin{tabular}{lrrrr}
\toprule
Mode & $k$ & Speedup & Latency (ms) & MSE \\
\midrule
Draft         & 0  & \textbf{96.9$\times$} & 0.152  & 0.169 \\
Partial-AR    & 1  & 48.5$\times$          & 0.304  & 0.169 \\
Partial-AR    & 4  & 19.4$\times$          & 0.760  & 0.169 \\
Partial-AR    & 16 & 5.7$\times$           & 2.586  & 0.167 \\
Full-AR       & 96 & 1.0$\times$           & 14{,}754 & 0.241 \\
\bottomrule
\end{tabular}
\end{table}

Table~\ref{tab:exp4} shows that draft mode achieves a \textbf{96.9$\times$} wall-clock speedup at MSE 0.169.
A single AR refinement step ($k=1$, 48.5$\times$ speedup) matches draft quality exactly, making additional AR steps above $k=1$ provide no accuracy benefit at greatly increased cost.

Full-AR mode ($k=96$) yields the \emph{worst} MSE of 0.241 — 42\% higher than draft.
This counter-intuitive result follows directly from the TiDAR training objective: the backbone is trained to predict all $H$ horizon positions in a single parallel forward pass, so sequential autoregressive decoding at test time is out-of-distribution relative to training.
The AR head functions as an \emph{auxiliary regulariser} that improves draft-head quality during training; it is not trained to support standalone sequential inference.
This finding is consistent across datasets: ETTh1 H\,=\,96 yields an 80.8$\times$ draft speedup (MSE 0.640 vs.\ full-AR 0.776), confirming that parallel draft mode is the correct primary inference mode for TiDAR-TS.

\subsection{Experiment 4: Cross-Dataset Architecture Transfer}
\label{sec:generalization}

We test whether genomes discovered on one dataset transfer without re-running evolutionary search.
Two source genomes are evaluated: the best H\,=\,96 genome found on ETTh1 (Experiment~1, S-Mamba), and the best H\,=\,96 genome found on Exchange-Rate.
Each genome is trained from scratch on the target dataset and evaluated at matching horizons.
The baseline is an all-Mamba (Rec-1) genome of comparable depth.

\begin{table}[t]
\caption{Transfer from ETTh1 (S-Mamba, H=96 search genome). Test MSE of the transferred genome vs.\ the all-Mamba baseline on four target datasets. Positive $\Delta$ means baseline is worse (transfer helps).}
\label{tab:exp5_from_etth}
\centering
\small
\setlength{\tabcolsep}{4pt}
\begin{tabular}{lrrrrr}
\toprule
Target & $H$ & Searched & Baseline & $\Delta\%$ \\
\midrule
\multirow{3}{*}{ETTh2}
 & 96  & \textbf{0.444} & 0.649 & $+$31.6\% \\
 & 336 & \textbf{0.735} & 1.071 & $+$31.3\% \\
 & 720 & 2.173 & \textbf{1.365} & $-$59.2\% \\
\midrule
\multirow{3}{*}{ETTm1}
 & 96  & 0.601 & \textbf{0.551} & $-$9.2\% \\
 & 336 & 0.781 & \textbf{0.688} & $-$13.5\% \\
 & 720 & 0.779 & \textbf{0.738} & $-$5.7\% \\
\midrule
\multirow{3}{*}{ETTm2}
 & 96  & 0.273 & \textbf{0.251} & $-$8.9\% \\
 & 336 & \textbf{0.618} & 0.717 & $+$13.8\% \\
 & 720 & \textbf{0.769} & 1.210 & $+$36.5\% \\
\midrule
\multirow{3}{*}{Exchange}
 & 96  & \textbf{0.147} & 0.405 & $+$63.6\% \\
 & 336 & \textbf{0.718} & 0.825 & $+$13.0\% \\
 & 720 & \textbf{1.301} & 1.482 & $+$12.2\% \\
\bottomrule
\end{tabular}
\end{table}

Table~\ref{tab:exp5_from_etth} shows that the ETTh1 genome transfers decisively to temporally compatible targets: ETTh2 ($+$31\% at H\,=\,96 and 336) and Exchange-Rate ($+$63.6\% at H\,=\,96), which share hourly resolution and smooth trend dynamics.
Transfer to ETTm (15-minute resolution) is mixed — the genome underperforms at short horizons but recovers at H\,=\,720 — reflecting a boundary condition where the $4{\times}$ higher sampling rate introduces frequency content outside the hourly inductive bias.
The reverse direction (Exchange$\to$ETT) fails predictably: the Exchange-Rate genome matches the all-Mamba baseline on ETTh1/ETTm1 but underperforms by 28--33\% on ETTh2/ETTm2, confirming that transfer is bounded by source-target temporal compatibility rather than limited by the evolutionary search itself.

